\documentclass{article}
\usepackage{xcolor}
\usepackage{expl3}
\usepackage{tikz}

\ExplSyntaxOn

% ============================================
% DÉFINITION DU SYSTÈME DE PALETTES
% ============================================

% Créer une nouvelle palette
\cs_new_protected:Nn \palette_new:n
  {
    \seq_new:c { g_palette_#1_seq }
  }

% Ajouter une couleur RGB à une palette
\cs_new_protected:Nn \palette_add_color:nn
  {
    \seq_gput_right:cn { g_palette_#1_seq } { #2 }
    % Définir automatiquement la couleur avec un nom standard
    \int_set:Nn \l_tmpa_int { \seq_count:c { g_palette_#1_seq } }
    \definecolor{palette@#1@\int_use:N\l_tmpa_int}{rgb}{#2}
  }

% Obtenir le nombre de couleurs
\cs_new:Nn \palette_count:n
  {
    \seq_count:c { g_palette_#1_seq }
  }

% Obtenir une couleur par index
\cs_new:Nn \palette_get_color:nn
  {
    \seq_item:cn { g_palette_#1_seq } { #2 }
  }

% Obtenir le NOM d'une couleur (expandable, pour TikZ)
\cs_new:Nn \palette_get_color_name:nn
  {
    palette@#1@#2
  }

% Définir une couleur xcolor depuis une palette
\cs_new_protected:Nn \palette_define_xcolor:nnn
  {
    \definecolor{#1}{rgb}{\seq_item:cn{g_palette_#2_seq}{#3}}
  }

% Utiliser directement une couleur
\cs_new_protected:Nn \palette_use_color:nn
  {
    \color{palette@#1@#2}
  }

% Boucle sur toutes les couleurs d'une palette
\cs_new_protected:Nn \palette_map:nn
  {
    \int_step_inline:nn { \palette_count:n { #1 } }
      {
        \group_begin:
        \color{palette@#1@##1}
        #2
        \group_end:
      }
  }

% ============================================
% INTERFACE UTILISATEUR
% ============================================

\NewDocumentCommand \paletteNew { m }
  { \palette_new:n { #1 } }

\NewDocumentCommand \paletteAddColor { m m }
  { \palette_add_color:nn { #1 } { #2 } }

\NewDocumentCommand \paletteCount { m }
  { \palette_count:n { #1 } }

\NewDocumentCommand \paletteGetColor { m m }
  { \palette_get_color:nn { #1 } { #2 } }

\NewDocumentCommand \paletteDefineColor { m m m }
  { \palette_define_xcolor:nnn { #1 } { #2 } { #3 } }

\NewDocumentCommand \paletteColor { m m }
  { \palette_use_color:nn { #1 } { #2 } }

\NewDocumentCommand \paletteMap { m m }
  { \palette_map:nn { #1 } { #2 } }

\NewExpandableDocumentCommand \paletteColorName { m m }
  { \palette_get_color_name:nn { #1 } { #2 } }

\ExplSyntaxOff

% ============================================
% EXEMPLE D'UTILISATION
% ============================================

\begin{document}

% Créer palette "Blackbody"
\paletteNew{Blackbody}
\paletteAddColor{Blackbody}{0.0, 0.0, 0.0}
\paletteAddColor{Blackbody}{0.4, 0.0, 0.2}
\paletteAddColor{Blackbody}{0.8, 0.2, 0.0}
\paletteAddColor{Blackbody}{1.0, 0.6, 0.0}
\paletteAddColor{Blackbody}{1.0, 1.0, 0.4}
\paletteAddColor{Blackbody}{1.0, 1.0, 1.0}

% Créer palette "Ocean"
\paletteNew{Ocean}
\paletteAddColor{Ocean}{0.0, 0.1, 0.3}
\paletteAddColor{Ocean}{0.0, 0.3, 0.6}
\paletteAddColor{Ocean}{0.2, 0.5, 0.8}
\paletteAddColor{Ocean}{0.5, 0.7, 0.9}
\paletteAddColor{Ocean}{0.8, 0.9, 1.0}

\section*{Démonstration des palettes}

\subsection*{Palette Blackbody}
Nombre de couleurs : \paletteCount{Blackbody}

\bigskip

% Afficher toutes les couleurs
\paletteMap{Blackbody}{\rule{2cm}{1cm}\quad}

\bigskip

% Utilisation dans du texte
{\paletteColor{Blackbody}{1} Noir} \quad
{\paletteColor{Blackbody}{3} Rouge-orangé} \quad
{\paletteColor{Blackbody}{5} Jaune pâle} \quad
{\paletteColor{Blackbody}{6} Blanc}

\subsection*{Palette Ocean}
Nombre de couleurs : \paletteCount{Ocean}

\bigskip

\paletteMap{Ocean}{\rule{2cm}{1cm}\quad}

\subsection*{Définir des couleurs nommées}

% Créer des couleurs xcolor depuis la palette
\paletteDefineColor{myred}{Blackbody}{3}
\paletteDefineColor{myblue}{Ocean}{3}

\textcolor{myred}{Texte en rouge} et \textcolor{myblue}{texte en bleu}.

\subsection*{Accès direct aux valeurs RGB}

La couleur 2 de Blackbody a les valeurs RGB : \paletteGetColor{Blackbody}{2}

\subsection*{Utilisation avec TikZ}

% Méthode 1 : Définir des couleurs nommées personnalisées
\paletteDefineColor{bb1}{Blackbody}{1}
\paletteDefineColor{bb2}{Blackbody}{2}
\paletteDefineColor{bb3}{Blackbody}{3}
\paletteDefineColor{bb4}{Blackbody}{4}
\paletteDefineColor{bb5}{Blackbody}{5}
\paletteDefineColor{bb6}{Blackbody}{6}

\bigskip

\begin{tikzpicture}
  % Méthode 1 : Utiliser les noms personnalisés
  \foreach \i in {1,...,6} {
    \fill[bb\i] (\i-1,0) rectangle (\i,1);
  }
  \node at (3,-0.5) {Méthode 1 : Noms personnalisés};
  
  % Méthode 2 : Utiliser directement \paletteColorName (EXPANDABLE!)
  \foreach \i in {1,...,5} {
    \fill[\paletteColorName{Ocean}{\i}] (\i-1,2) rectangle (\i,3);
  }
  \node at (2.5,1.5) {Méthode 2 : \textbackslash paletteColorName};
  
  % Cercles avec Blackbody
  \foreach \i in {1,...,6} {
    \fill[\paletteColorName{Blackbody}{\i}] (7+\i*0.7, 0.5) circle (0.3);
  }
  
  % Cercles avec Ocean
  \foreach \i in {1,...,5} {
    \fill[\paletteColorName{Ocean}{\i}] (7+\i*0.7, 2.5) circle (0.3);
  }
\end{tikzpicture}

\bigskip

\textbf{Deux façons d'utiliser les couleurs :}
\begin{itemize}
\item \texttt{\textbackslash paletteDefineColor\{mycolor\}\{Palette\}\{index\}} : Crée un nom personnalisé réutilisable
\item \texttt{\textbackslash paletteColorName\{Palette\}\{index\}} : Retourne directement le nom de couleur (expandable, utilisable dans TikZ !)
\end{itemize}

Exemple : \texttt{\textbackslash fill[\textbackslash paletteColorName\{Ocean\}\{3\}] ...} fonctionne directement !

\end{document}